% !TeX root = ../MQF_Cap_13_Python_Asset_Allocation_ALM.tex
% Capitoli/MQF_Cap_13_Python_Asset_Allocation_ALM.tex


\chapter[Python e Asset--Liability Management]{Python e Asset--Liability Management: ottimizzazione, sensitivit\`a e scenari}

\label{cap:python-alm}

I Capitoli 11 e 12 hanno costruito il quadro teorico necessario per
formulare e interpretare problemi di allocazione e di Asset--Liability Management mediante programmazione lineare. Il Capitolo 11 ha
introdotto la struttura generale del problema lineare, i vincoli attivi, la
dualit\`a e l'analisi di sensitivit\`a; il Capitolo 12 ha esteso questa
impostazione a un problema ALM deterministico multi-periodale, nel quale
disponibilit\`a, fabbisogni, giacenze e funding devono essere coordinati lungo
una sequenza di date.

Questo capitolo compie il passaggio successivo. L'obiettivo non \`e
introdurre una nuova classe di modelli, ma rendere operativo in Python il
problema ALM gi\`a formulato, traducendolo nella forma richiesta dal solver e
utilizzandolo per analizzare configurazioni alternative.

Il filo conduttore \`e un Caso Aula ispirato alla crisi di Silicon Valley
Bank del 2023. La configurazione di riferimento viene dapprima risolta e
verificata come benchmark, ricostruendo composizione dell'attivo, giacenze di
liquidit\`a, ricorso al funding, vincoli attivi e valori marginali.

Il capitolo introduce inoltre un piccolo focus Python dedicato alle funzioni
riutilizzabili. Questo strumento verr\`a impiegato per risolvere lo stesso
modello al variare di tre driver: il parametro $\theta$, che modifica i
fabbisogni iniziali di liquidit\`a; il parametro $\phi$, che modifica la
capacit\`a di funding; il parametro $\psi$, che modifica il costo del funding
e la corrispondente struttura di rimborso.

L'analisi distingue quindi tre problemi. La sensitivit\`a studia come cambia
la soluzione quando varia un driver alla volta; la frontiera di
fattibilit\`a individua fino a quale livello di stress esiste ancora un piano
ammissibile; gli scenari combinano simultaneamente pi\`u variazioni delle
condizioni del problema.

Particolare attenzione viene dedicata ai controlli. Una soluzione dichiarata
ottima dal solver deve essere ricondotta alla formulazione economica
originaria verificando vincoli, residui, slack, bound attivi e significato
dei valori marginali.

Il metodo di lavoro rimane quello introdotto nel Capitolo 4 per i capitoli
applicativi e viene qui assunto come acquisito.

Accanto al Caso Aula, un Caso TakeHome ispirato alla crisi di First Republic
Bank del 2023 trasferisce la stessa architettura di analisi a una diversa
struttura dell'attivo, a due ondate di deposit run e a ulteriori vincoli
strutturali.

Il percorso complessivo del capitolo \`e quindi

$$
\begin{array}{c}
\text{modello ALM}
\\[3pt]
\downarrow
\\[3pt]
\text{forma computazionale}
\\[3pt]
\downarrow
\\[3pt]
\text{benchmark verificato}
\\[3pt]
\downarrow
\\[3pt]
\text{sensitivit\`a, fattibilit\`a e scenari}.
\end{array}
$$

Python rende ripetibile la soluzione del problema; il modello economico
rimane il riferimento necessario per interpretarne i risultati.

\section{Python per costruire funzioni riutilizzabili}

Nei capitoli applicativi precedenti, Python \`e stato utilizzato per
organizzare traiettorie, operare su matrici e selezionare scenari. Nel
presente capitolo, emerge una necessit\`a differente: eseguire ripetutamente la
stessa procedura di ottimizzazione modificandone alcuni input.

Una funzione permette di racchiudere una sequenza di operazioni in una
struttura riutilizzabile. In Python, viene definita mediante la parola chiave
\texttt{def}:

\begin{verbatim}
def funzione(input):
    elaborazione
    return output
\end{verbatim}

Gli argomenti rappresentano le informazioni ricevute dalla funzione, mentre
\texttt{return} specifica i risultati restituiti al termine
dell'elaborazione.

Ad esempio,

\begin{verbatim}
def scala(x, a):
    y = a * x
    return y
\end{verbatim}

\noindent applica la stessa operazione a valori differenti degli input senza
dover riscrivere le istruzioni che la realizzano.

Nel problema ALM, la funzione dovr\`a ricevere i parametri che identificano
una particolare configurazione, costruire gli oggetti necessari al solver ed
estrarre dalla soluzione le quantit\`a rilevanti. La struttura sar\`a quindi
del tipo

$$
(\theta,\phi,\psi)
\longrightarrow
\texttt{solve\_alm}(\theta,\phi,\psi)
\longrightarrow
(\Pi^*,x^*,g^*,u^*,\ldots).
$$

La costruzione della funzione richiede di distinguere gli elementi che
rimangono invariati tra le diverse configurazioni da quelli che devono essere
aggiornati quando cambiano gli input. Se un parametro interviene soltanto su
una componente della formulazione, tutte le altre devono rimanere
inalterate.

La funzione costituisce quindi uno strumento per rendere ripetibile e
controllabile la procedura di soluzione. Prima di costruirla nel caso
specifico, occorre tuttavia tradurre con precisione il modello ALM nella forma
richiesta dal solver.


\section{Dal modello ALM al problema computazionale}

Il Capitolo 12 ha formulato il problema di Asset--Liability Management
deterministico come un programma lineare nel quale le decisioni di
allocazione iniziale devono essere compatibili con una sequenza di fabbisogni
di liquidit\`a, con la possibilit\`a di mantenere giacenze tra date successive
e con un ricorso limitato al funding esterno.

Per risolvere numericamente il problema, questa formulazione deve essere
tradotta negli oggetti richiesti dal solver. Il passaggio non introduce un
nuovo modello: modifica la rappresentazione dello stesso problema, passando
dalla forma economico-matematica alla forma matriciale utilizzata
dall'algoritmo.

Con \texttt{scipy.optimize.linprog}, il problema viene espresso nella forma

\[
\min_{v}\; f^{\prime}v
\]

soggetto a

\[
A_{\mathrm{ub}}v\leq b_{\mathrm{ub}},
\]

\[
A_{\mathrm{eq}}v=b_{\mathrm{eq}},
\]

\noindent e ai bound assegnati alle singole variabili.

La traduzione richiede quindi quattro scelte coordinate: l'ordinamento delle
variabili decisionali, la costruzione del vettore della funzione obiettivo,
la rappresentazione matriciale dei vincoli e la definizione dei bound.

\subsection{Ordinare le variabili decisionali}

Nel Caso Aula, il modello comprende quattro variabili di allocazione iniziale,

\[
x_1,x_2,x_3,x_4,
\]

\noindent quattro giacenze di liquidit\`a,

\[
g_1,g_2,g_3,g_4,
\]

\noindent e tre variabili di funding,

\[
u_1,u_2,u_3.
\]

Il vettore delle variabili decisionali viene quindi ordinato come

\[
v=
\left(
x_1,x_2,x_3,x_4,
g_1,g_2,g_3,g_4,
u_1,u_2,u_3
\right)^{\prime}.
\]

L'ordine non ha significato economico autonomo, ma deve rimanere invariato in
tutti gli oggetti costruiti successivamente. Ogni coefficiente della funzione
obiettivo, ogni colonna delle matrici dei vincoli e ogni bound devono riferirsi
alla stessa posizione del vettore $v$.

Questa coerenza \`e essenziale soprattutto nei problemi multi-periodali, nei
quali variabili economicamente diverse compaiono contemporaneamente nella
stessa equazione di bilancio.

\subsection{Tradurre la funzione obiettivo}

Nel modello ALM, la funzione obiettivo massimizza il rendimento netto del
piano. Con l'ordinamento adottato per $v$, il vettore dei coefficienti
economici pu\`o essere scritto come

\[
c_{\Pi}
=
\left(
c_1,c_2,c_3,c_4,
r_1^g,r_2^g,r_3^g,0,
-r_1^u,-r_2^u,-r_3^u
\right)^{\prime}.
\]

La funzione obiettivo pu\`o essere scritta come

\[
\Pi(v)
=
c_{\Pi}^{\prime}v.
\]

\noindent Il problema richiede quindi di risolvere

\[
\max_{v}\; c_{\Pi}^{\prime}v.
\]

Il coefficiente associato a $g_4$ \`e nullo perch\'e la giacenza terminale non
genera un ulteriore rendimento all'interno dell'orizzonte considerato. I
coefficienti delle variabili $u_t$ sono invece negativi, poich\'e il funding
comporta un costo.

La funzione \texttt{linprog} risolve problemi di minimo. Per utilizzare lo
stesso modello senza modificarne il significato economico, si definisce

\[
f=-c_{\Pi}
\]

\noindent e si risolve

\[
\min_{v}\; f^{\prime}v.
\]

Poich\'e

\[
f^{\prime}v
=
-c_{\Pi}^{\prime}v
=
-\Pi(v),
\]

\noindent minimizzare $f^{\prime}v$ equivale a massimizzare $\Pi(v)$.

La distinzione tra $c_{\Pi}$ e $f$ deve rimanere esplicita. Il primo contiene
i coefficienti della funzione economica che si desidera massimizzare; il
secondo \`e il vettore effettivamente fornito al solver.

\subsection{Costruire i vincoli di uguaglianza}

I vincoli di uguaglianza comprendono il vincolo sulle risorse iniziali e i
bilanci di liquidit\`a alle diverse date.

Il vincolo iniziale \`e

\[
x_1+x_2+x_3+x_4=A_0.
\]

I bilanci temporali seguono invece la struttura introdotta nel Capitolo 12.
Alla prima data,

\[
a_1^{\prime}x+u_1-g_1=d_1.
\]

Per le date intermedie,

\[
a_t^{\prime}x
+
(1+r_{t-1}^{g})g_{t-1}
+
u_t
-
(1+r_{t-1}^{u})u_{t-1}
-
g_t
=
d_t,
\]

\noindent mentre, alla data terminale,

\[
a_4^{\prime}x
+
(1+r_3^{g})g_3
-
(1+r_3^{u})u_3
-
g_4
=
d_4.
\]

Queste equazioni vengono raccolte nella forma

\[
A_{\mathrm{eq}}v=b_{\mathrm{eq}}.
\]

Nel Caso Aula, $A_{\mathrm{eq}}$ contiene quindi una riga per il vincolo sulle
risorse iniziali e una riga per ciascun bilancio temporale.

La matrice dei cash flow delle attivit\`a,

\[
A^{CF}
=
\begin{pmatrix}
a_{11} & a_{12} & a_{13} & a_{14}\\
a_{21} & a_{22} & a_{23} & a_{24}\\
a_{31} & a_{32} & a_{33} & a_{34}\\
a_{41} & a_{42} & a_{43} & a_{44}
\end{pmatrix},
\]

\noindent determina i coefficienti delle variabili $x_j$ nelle quattro equazioni
temporali. Le giacenze e il funding compaiono invece con coefficienti che
riflettono il trasferimento di liquidit\`a e il rimborso tra una data e la
successiva.

Il vettore dei termini noti assume la forma

\[
b_{\mathrm{eq}}
=
\left(
A_0,
d_1,
d_2,
d_3,
d_4
\right)^{\prime}.
\]

La costruzione di $A_{\mathrm{eq}}$ rende quindi esplicita la dimensione
temporale del problema: una variabile scelta o generata a una data pu\`o
entrare nei bilanci delle date successive attraverso i coefficienti di
riporto o di rimborso.

\subsection{Disuguaglianze e bound}

I vincoli prudenziali che non richiedono un'uguaglianza vengono raccolti
nella forma

\[
A_{\mathrm{ub}}v\leq b_{\mathrm{ub}}.
\]

Nel Caso Aula, rientrano in questa struttura il limite alla perdita di stress
del portafoglio e il vincolo di concentrazione sull'attivit\`a alla quale
viene imposto un limite massimo.

Il vincolo di stress,

\[
\ell^{\prime}x\leq K,
\]

\noindent dipende soltanto dalle variabili di allocazione. Le colonne corrispondenti a
$g$ e $u$ hanno quindi coefficiente nullo.

Analogamente, un limite

\[
x_3\leq \overline{x}_3
\]

\noindent viene rappresentato da una riga nella quale compare un unico coefficiente
unitario nella posizione associata a $x_3$.

Non tutti i limiti sulle variabili devono tuttavia essere trasformati in
righe di $A_{\mathrm{ub}}$. Quando il vincolo riguarda direttamente
l'intervallo ammissibile di una singola variabile, \texttt{linprog} permette
di rappresentarlo mediante i bound.

Per le allocazioni e le giacenze, si impone

\[
x_j\geq 0,
\qquad
g_t\geq 0.
\]

Per il funding,

\[
0\leq u_t\leq \overline{u}_t.
\]

I bound svolgono quindi una funzione distinta dalle disuguaglianze generali:
definiscono direttamente l'intervallo ammissibile di ciascuna componente di
$v$.

\subsection{Dalla rappresentazione matriciale al solver}

Una volta costruiti gli oggetti

\[
f,
\qquad
A_{\mathrm{ub}},
\qquad
b_{\mathrm{ub}},
\qquad
A_{\mathrm{eq}},
\qquad
b_{\mathrm{eq}},
\qquad
\text{bound},
\]

\noindent il problema pu\`o essere passato al solver.

In Python, la chiamata assume la forma essenziale

\begin{verbatim}
res = linprog(
    c=f,
    A_ub=A_ub,
    b_ub=b_ub,
    A_eq=A_eq,
    b_eq=b_eq,
    bounds=bound,
    method="highs"
)
\end{verbatim}


Nella chiamata a \texttt{linprog}, i nomi che compaiono a sinistra del segno
\texttt{=} sono gli argomenti previsti dalla funzione, mentre quelli che
compaiono a destra sono le variabili costruite nel notebook.

Quando i due nomi coincidono, come in {\verb|A_eq=A_eq|}, non si tratta di un'uguaglianza matematica, ma dell'assegnazione della
variabile \texttt{A\_eq} del notebook all'argomento \texttt{A\_eq} richiesto
da \texttt{linprog}. La coincidenza dei nomi dipende soltanto dalla
convenzione adottata nel codice.

Il caso {\verb|c=f|} \`e invece diverso. La funzione \texttt{linprog} utilizza il nome
\texttt{c} per il vettore dei coefficienti della funzione obiettivo da
minimizzare. Nel notebook, lo stesso oggetto \`e stato indicato con $f$,
dove $f=-c_{\Pi}$, per trasformare il problema originario di massimo nel problema equivalente
di minimo.

La scrittura {\verb|c=f|} significa quindi che all'argomento \texttt{c} previsto dal solver viene
passato il vettore $f$ costruito nel notebook.


L'oggetto \texttt{res} non contiene soltanto il vettore della soluzione
ottima. Raccoglie anche lo stato della procedura, il valore della funzione
obiettivo, residui, slack e informazioni marginali associate ai vincoli.

La soluzione restituita dal solver deve inoltre essere ricondotta alla
convenzione economica originaria. Poich\'e il problema numerico minimizza
$-\Pi(v)$, il valore economico ottimo viene ricostruito come

\[
\Pi^*
=
-\texttt{res.fun}.
\]

Analogamente, le componenti di \texttt{res.x} devono essere ricondotte ai
blocchi

\[
x^*,
\qquad
g^*,
\qquad
u^*,
\]

\noindent secondo l'ordinamento scelto inizialmente.

La traduzione computazionale pu\`o quindi essere sintetizzata come

\[
\begin{array}{ccc}
\text{modello economico}
&
\longrightarrow
&
\text{oggetti del solver}
\\[4pt]
\Pi(v)
&
\longrightarrow
&
f
\\[4pt]
\text{bilanci temporali}
&
\longrightarrow
&
A_{\mathrm{eq}},\,b_{\mathrm{eq}}
\\[4pt]
\text{vincoli prudenziali}
&
\longrightarrow
&
A_{\mathrm{ub}},\,b_{\mathrm{ub}}
\\[4pt]
\text{limiti sulle variabili}
&
\longrightarrow
&
\text{bound}.
\end{array}
\]

La corrispondenza deve essere verificata prima di interpretare qualsiasi
risultato numerico. Una soluzione ottenuta da un problema codificato in modo
errato pu\`o essere perfettamente ottima per il solver e al tempo stesso non
rappresentare il problema ALM che si intende analizzare.

Nel seguito, questa rappresentazione viene applicata alla configurazione
benchmark del Caso Aula. La prima verifica riguarder\`a quindi non soltanto
l'esistenza di una soluzione ottima, ma la coerenza tra tale soluzione, i
vincoli del modello e il significato economico del piano ALM.


\section{Dal benchmark al modello parametrico}

La configurazione benchmark fornisce una soluzione di riferimento. Il
passaggio successivo consiste nel modificare alcune condizioni economiche e
osservare come reagiscano il piano ALM e il suo valore.

Nel Caso Aula, vengono introdotti tre driver,

$$
\theta,
\qquad
\phi,
\qquad
\psi,
$$

\noindent che agiscono rispettivamente sui fabbisogni di liquidit\`a, sulla capacit\`a
di funding e sul costo del funding.

\subsection{Tre driver con effetti differenti}

Il parametro $\theta$ modifica i fabbisogni delle prime due date:

$$
d(\theta)
=
\left(
30\theta,
30\theta,
22,
15
\right)^{\prime}.
$$

Il benchmark corrisponde a $\theta=1$, mentre valori superiori a uno
rappresentano un aumento dei fabbisogni iniziali.

Il parametro $\phi$ modifica la capacit\`a massima di funding. Se

$$
\overline{u}
=
\left(
10,8,6
\right)^{\prime},
$$

\noindent si pone

$$
\overline{u}(\phi)
=
\phi\overline{u}.
$$

Valori di $\phi<1$ rappresentano quindi una contrazione della capacit\`a
disponibile.

Il parametro $\psi$ modifica invece il costo del funding:

$$
r^u(\psi)
=
\psi r^u.
$$

La sua variazione interessa sia la funzione obiettivo sia i bilanci
temporali, poich\'e il funding ottenuto a una data deve essere rimborsato alla
data successiva con il relativo costo.

\subsection{Dove entrano i parametri nel solver}

I tre driver modificano componenti differenti della formulazione.

Poich\'e $\theta$ interviene nei fabbisogni $d_1$ e $d_2$, modifica il vettore
dei termini noti:

$$
b_{\mathrm{eq}}(\theta)
=
\left(
A_0,
30\theta,
30\theta,
22,
15
\right)^{\prime}.
$$

Il parametro $\phi$ modifica invece i bound delle variabili di funding:

$$
0
\leq
u_t
\leq
\phi\overline{u}_t.
$$

Infine, $\psi$ modifica i coefficienti del funding nella funzione obiettivo e
i coefficienti di rimborso nelle equazioni temporali:

$$
r_t^u(\psi)
=
\psi r_t^u,
$$

$$
1+r_t^u(\psi)
=
1+\psi r_t^u.
$$

La corrispondenza pu\`o quindi essere sintetizzata come

$$
\begin{array}{ccl}
\theta
&\longrightarrow&
b_{\mathrm{eq}},
\\[4pt]
\phi
&\longrightarrow&
\text{bound},
\\[4pt]
\psi
&\longrightarrow&
f,\;A_{\mathrm{eq}}.
\end{array}
$$

Questa distinzione \`e essenziale: la variazione di un parametro deve essere
trasferita esclusivamente agli oggetti della formulazione che ne dipendono.

\subsection{La funzione di soluzione}

Le dipendenze precedenti vengono incorporate nella funzione

\begin{verbatim}
def solve_alm(theta=1.0, phi=1.0, psi=1.0):
    ...
    return ...
\end{verbatim}

I valori predefiniti $\theta=\phi=\psi=1$ identificano il benchmark. Per ogni
altra configurazione, la funzione aggiorna gli oggetti dipendenti dai
parametri, risolve il problema e restituisce almeno

$$
\Pi^*,
\qquad
x^*,
\qquad
g^*,
\qquad
u^*,
$$

\noindent insieme all'informazione sulla fattibilit\`a e, quando necessario, agli
elementi utilizzati per la diagnostica della soluzione.

La stessa procedura pu\`o quindi essere richiamata su una griglia di valori
senza modificare manualmente la formulazione del problema. L'analisi pu\`o
cos\`i concentrarsi sul confronto tra le soluzioni e sui cambiamenti dei
vincoli attivi, anzich\'e sulla ripetizione del codice.

Su questa base, vengono costruite le analisi di sensitivit\`a, la ricerca
della frontiera di fattibilit\`a e il confronto tra scenari di stress.

\section{Sensitivit\`a, fattibilit\`a e scenari di stress}

La funzione parametrica permette di risolvere ripetutamente lo stesso modello
ALM al variare dei driver $\theta$, $\phi$ e $\psi$. L'analisi pu\`o quindi
spostarsi dalla singola configurazione benchmark al comportamento della
soluzione in un insieme di condizioni alternative.

Tre domande devono essere mantenute distinte. La prima riguarda la
\emph{sensitivit\`a}: come cambiano il valore ottimo e il piano ALM quando
varia un parametro alla volta? La seconda riguarda la
\emph{fattibilit\`a}: fino a quale intensit\`a dello stress esiste ancora una
soluzione che soddisfa tutti i vincoli? La terza riguarda gli
\emph{scenari}: che cosa accade quando pi\`u condizioni avverse vengono
modificate simultaneamente?

\subsection{Sensitivit\`a rispetto ai fabbisogni di liquidit\`a}

Il parametro $\theta$ modifica i fabbisogni delle prime due date secondo

\[
d(\theta)
=
\left(
30\theta,
30\theta,
22,
15
\right)^{\prime}.
\]

L'analisi considera una griglia di valori di $\theta$ mantenendo

\[
\phi=1,
\qquad
\psi=1.
\]

Per ciascuna configurazione, vengono osservati almeno

\[
\Pi^*(\theta),
\qquad
x^*(\theta),
\qquad
g^*(\theta),
\qquad
u^*(\theta).
\]

Un aumento di $\theta$ rende pi\`u oneroso il soddisfacimento dei fabbisogni
iniziali. Il valore ottimo tende quindi a ridursi e il piano ALM deve
riorganizzare l'impiego delle attivit\`a, delle giacenze e del funding per
mantenere soddisfatti i bilanci temporali.

Nel benchmark, la capacit\`a di funding alla terza data \`e gi\`a interamente
utilizzata,

\[
u_3^*=\overline{u}_3.
\]

L'aumento dei fabbisogni iniziali opera quindi in un sistema nel quale una
delle risorse di liquidit\`a future \`e gi\`a al proprio limite. La
sensitivit\`a rispetto a $\theta$ deve essere letta alla luce di questa
struttura: l'effetto di uno stress non dipende soltanto dalla sua ampiezza,
ma anche dai vincoli che risultano attivi nella configurazione di partenza.

L'andamento di

\[
\theta
\longmapsto
\Pi^*(\theta)
\]

\noindent non deve inoltre essere interpretato come necessariamente caratterizzato da
un'unica relazione marginale valida sull'intero intervallo. Finch\'e rimane
invariata la struttura dei vincoli che determinano l'ottimo, la funzione
valore pu\`o seguire uno stesso regime locale. Quando cambia l'insieme dei
vincoli attivi, cambia anche la risposta marginale del problema.

\subsection{Sensitivit\`a rispetto alla capacit\`a di funding}

Il parametro $\phi$ modifica i limiti superiori del funding:

\[
0
\leq
u_t
\leq
\phi\overline{u}_t.
\]

In questo caso, lo stress non modifica direttamente i fabbisogni, ma la
quantit\`a massima di risorse esterne che pu\`o essere utilizzata per
soddisfarli.

L'analisi di

\[
\phi
\longmapsto
\Pi^*(1,\phi,1)
\]

\noindent permette quindi di distinguere due situazioni economicamente differenti.
Quando una capacit\`a di funding \`e scarsa e il relativo bound \`e attivo,
una variazione di $\phi$ pu\`o modificare il piano ottimo e il valore
dell'obiettivo. Quando invece la capacit\`a disponibile eccede quella che il
piano desidera utilizzare, ulteriori incrementi di $\phi$ non producono
necessariamente alcun beneficio.

Pu\`o quindi emergere una regione nella quale

\[
u_t^*<\phi\overline{u}_t,
\]

\noindent e il corrispondente bound superiore non \`e attivo. In tale regione, un ulteriore aumento della capacit\`a di funding non modifica localmente la soluzione ottima.

Questa osservazione riprende in forma numerica il significato dei prezzi
ombra sviluppato nei Capitoli 11 e 12: una risorsa aggiuntiva ha valore
marginale soltanto quando il vincolo che la limita \`e effettivamente
rilevante per l'ottimo.

\subsection{Sensitivit\`a rispetto al costo del funding}

Il parametro $\psi$ agisce sul vettore

\[
r^u(\psi)
=
\psi r^u.
\]

A differenza di $\theta$ e $\phi$, la sua variazione modifica
simultaneamente due componenti del problema. Aumentano i costi del funding
nella funzione obiettivo e aumentano i coefficienti di rimborso che compaiono
nei bilanci delle date successive.

L'analisi

\[
\psi
\longmapsto
\Pi^*(1,1,\psi)
\]

\noindent deve quindi essere interpretata tenendo conto di entrambi gli effetti.

Un aumento di $\psi$ rende il funding pi\`u costoso, ma non implica necessariamente una riduzione del suo utilizzo. Se il funding \`e necessario per soddisfare i bilanci temporali, il modello pu\`o continuare a ricorrervi anche a condizioni meno favorevoli. L'aggiustamento pu\`o allora avvenire, quando possibile, attraverso una diversa composizione dell'attivo, delle giacenze, o delle altre fonti di liquidit\`a.

La sensitivit\`a rispetto a $\psi$ consente quindi di distinguere due effetti: da un lato, il contributo del funding alla fattibilit\`a del piano; dall'altro, il suo impatto sul valore economico della soluzione.

\subsection{Cambiamenti di regime e vincoli attivi}

Le tre analisi precedenti non devono essere ridotte al confronto tra valori
successivi di $\Pi^*$.

Per ogni parametro, \`e utile osservare anche se cambiano

\[
x^*,
\qquad
g^*,
\qquad
u^*,
\]

\noindent e soprattutto se cambia l'insieme dei vincoli attivi.

In un programma lineare parametrico, un intervallo di valori pu\`o essere
caratterizzato dalla stessa struttura dell'ottimo. Entro tale intervallo,
piccole variazioni del parametro producono una risposta coerente con i
marginali associati alla configurazione corrente.

Quando un bound viene raggiunto, un vincolo precedentemente attivo cessa di
esserlo oppure un nuovo vincolo diventa vincolante, la struttura della
soluzione cambia. I valori marginali della configurazione precedente non
devono allora essere estesi automaticamente oltre quel punto.

L'analisi numerica su griglia permette di rendere visibili questi cambiamenti
e di collegare la sensitivit\`a locale studiata teoricamente alla variazione
della soluzione su intervalli pi\`u ampi.

\subsection{Dalla sensitivit\`a alla frontiera di fattibilit\`a}

L'aumento di $\theta$ pone anche una domanda diversa da quella precedente.
Non interessa pi\`u soltanto sapere come cambia l'ottimo, ma determinare fino
a quale livello di fabbisogno il sistema sia ancora in grado di soddisfare
simultaneamente tutti i vincoli.

Si passa quindi dalla domanda

\[
\text{come cambia la soluzione ottima?}
\]

\noindent alla domanda

\[
\text{esiste ancora una soluzione ammissibile?}
\]

Per una successione crescente di valori di $\theta$, la funzione parametrica
permette di verificare se la corrispondente regione ammissibile

$$
\mathcal{F}(\theta)
=
\left\{
v:
A_{\mathrm{eq}}v=b_{\mathrm{eq}}(\theta),
\;
A_{\mathrm{ub}}v\leq b_{\mathrm{ub}},
\;
v\in\text{bound}
\right\}
$$

\noindent sia non vuota oppure vuota.

Finch\'e

$$
\mathcal{F}(\theta)\neq\emptyset,
$$

\noindent esiste almeno una combinazione delle variabili decisionali compatibile con
tutti i vincoli del problema. All'aumentare di $\theta$, diventano progressivamente pi\`u stringenti i requisiti di liquidit\`a che una soluzione ammissibile deve soddisfare; oltre $\theta_{\mathrm{crit}}$, nessun vettore di decisione soddisfa simultaneamente tutti i vincoli. 
Definiamo quindi

$$
\theta_{\mathrm{crit}}
=
\max
\left\{
\theta:
\mathcal{F}(\theta)\neq\emptyset
\right\},
$$

\noindent ossia il massimo livello di stress compatibile con l'esistenza di un piano
ALM ammissibile.

Se una griglia individua due valori consecutivi $\theta_F<\theta_I$, con problema feasible in $\theta_F$ e infeasible in $\theta_I$, si pu\`o
localizzare la frontiera nell'intervallo

$$
\theta_{\mathrm{crit}}
\in
[\theta_F,\theta_I).
$$

Una griglia pi\`u fine consente successivamente di restringere tale
intervallo.

La frontiera di fattibilit\`a non rappresenta quindi soltanto il punto nel
quale il valore ottimo peggiora ulteriormente. Essa separa due situazioni
qualitativamente differenti. Per valori di $\theta$ inferiori o uguali alla
soglia critica, esistono valori delle variabili decisionali

$$
v=(x,g,u)
$$

\noindent che appartengono alla regione ammissibile e soddisfano simultaneamente tutti
i vincoli. Per valori superiori a $\theta_{\mathrm{crit}}$, invece,

$$
\mathcal{F}(\theta)=\emptyset,
$$

\noindent e non esiste alcuna combinazione di allocazioni, giacenze e funding capace di
soddisfare i bilanci temporali e gli altri vincoli del modello.

Il valore ottenuto mediante la ricerca su griglia rimane una localizzazione
numerica della frontiera e non deve essere interpretato come una soglia
analitica esatta.


\subsection{Dalle variazioni univariate agli scenari}

Le analisi di sensitivit\`a modificano un parametro alla volta. Questa scelta
permette di isolare il ruolo di ciascun driver, ma non rappresenta una
situazione nella quale pi\`u condizioni avverse si manifestano
contemporaneamente.

Per questo motivo, vengono considerate tre configurazioni con variazioni
congiunte:

\[
\begin{array}{c|ccc}
\text{Scenario}
&
\theta
&
\phi
&
\psi
\\ \hline
\text{Standard}
&
1.00
&
1.00
&
1.00
\\
\text{Mediamente critico}
&
1.03
&
0.90
&
1.25
\\
\text{Critico}
&
1.06
&
0.70
&
1.50
\end{array}
\]

Lo scenario Standard coincide con il benchmark. Negli altri due scenari,
aumentano i fabbisogni iniziali, si riduce la capacit\`a di funding e aumenta
il relativo costo.

Gli scenari non sono stati probabilizzati e non rappresentano quindi stati
aleatori ai quali associare una probabilit\`a di accadimento. Sono
configurazioni deterministiche utilizzate per confrontare la risposta del
modello a combinazioni differenti dei driver.

Il confronto non deve limitarsi a

\[
\Delta\Pi^*
=
\Pi^*_{\mathrm{scenario}}
-
\Pi^*_{\mathrm{Standard}}.
\]

Occorre osservare anche come cambiano la composizione dell'attivo, le
giacenze, il ricorso al funding e l'utilizzo delle capacit\`a disponibili.

Il passaggio dalla sensitivit\`a univariata agli scenari permette quindi di
distinguere due livelli di analisi:

\[
\begin{array}{ccc}
\text{sensitivit\`a}
&
&
\text{scenario}
\\[3pt]
\text{un driver alla volta}
&
\longrightarrow
&
\text{driver congiunti}.
\end{array}
\]

La prima consente di comprendere il ruolo specifico dei parametri e di
individuare cambiamenti nella struttura dell'ottimo; la seconda mostra come
gli stessi meccanismi interagiscano quando le condizioni del problema
peggiorano simultaneamente.

Sensitivit\`a, frontiera di fattibilit\`a e scenario analysis forniscono
quindi informazioni complementari. La soluzione benchmark descrive un piano;
l'analisi parametrica mostra quanto quel piano, e la struttura che lo
determina, siano robusti rispetto a variazioni delle condizioni economiche.



\section{Dal Caso Aula al Caso TakeHome: un diverso problema ALM}

Il Caso TakeHome trasferisce la stessa architettura di analisi a un problema
ALM con caratteristiche differenti. L'obiettivo non \`e replicare il Caso
Aula modificando soltanto alcuni dati, ma verificare se il metodo rimanga
applicabile quando cambiano la struttura dell'attivo, la distribuzione
temporale dei fabbisogni e alcuni vincoli del problema.

Il caso \`e ispirato alla crisi di First Republic Bank del 2023. Rispetto al
Caso Aula, il portafoglio comprende cinque classi di attivit\`a e presenta una
maggiore esposizione a componenti di bilancio meno liquide. La struttura dei
fabbisogni incorpora inoltre due differenti ondate di deposit run, mentre il
funding esterno rimane soggetto a limiti di capacit\`a.

La logica generale del problema rimane quella introdotta nel Capitolo 12:

\vspace{-0.3cm}

\begin{gather*}
\text{allocazione iniziale}\\
\downarrow\\
\text{cash flow temporali}\\
\downarrow\\
\text{giacenze e funding}\\
\downarrow\\
\text{rispetto dei fabbisogni}.
\end{gather*}


Cambiano per\`o gli oggetti attraverso i quali questa logica viene
implementata.

\subsection{Una diversa struttura dell'attivo}

Nel Caso TakeHome, il vettore delle allocazioni comprende cinque componenti,

$$
x=
(x_1,x_2,x_3,x_4,x_5)^{\prime},
$$

\noindent associate rispettivamente a liquidit\`a, titoli a breve termine, mutui
residenziali a lunga scadenza, prestiti immobiliari commerciali e altre forme
di credito alla clientela.

La maggiore articolazione dell'attivo modifica sia la matrice dei cash flow
sia la struttura dei vincoli prudenziali.

Accanto al limite massimo sull'esposizione verso una singola classe di
attivit\`a, compare infatti un vincolo legacy,

$$
x_3+x_4\geq 35,
$$

\noindent che impone di mantenere almeno una determinata quota del portafoglio nelle
componenti a pi\`u lunga scadenza.

Questo vincolo distingue in modo importante il Caso TakeHome dal Caso Aula.
La composizione dell'attivo non pu\`o infatti essere modificata liberamente
per aumentare la liquidit\`a disponibile: una parte del portafoglio rimane
vincolata a posizioni meno liquide.

\subsection{Due ondate di fabbisogno}

La parametrizzazione distingue due differenti componenti dello stress di
liquidit\`a.

Il parametro $\alpha$ modifica il fabbisogno alla prima data, mentre il
parametro $\beta$ modifica quello alla data terminale:

$$
d(\alpha,\beta)
=
\left(
32\alpha,
18,
14,
32\beta
\right)^{\prime}.
$$

La distinzione permette di confrontare due forme di stress che hanno la stessa
natura economica, ma si manifestano in momenti differenti dell'orizzonte.

Un aumento di $\alpha$ richiede maggiore liquidit\`a immediata e deve essere
assorbito utilizzando risorse disponibili nelle prime date. Un aumento di
$\beta$ agisce invece sul fabbisogno terminale e pu\`o essere assorbito, almeno
in parte, attraverso liquidit\`a accumulata nelle date precedenti.

Il confronto tra $\alpha$ e $\beta$ rende quindi nuovamente evidente che, in
un problema ALM multi-periodale, non conta soltanto l'ammontare complessivo
dei fabbisogni, ma anche il momento nel quale essi devono essere soddisfatti.

\subsection{Capacit\`a di funding e vincoli strutturali}

Anche nel Caso TakeHome, la capacit\`a di funding viene parametrizzata mediante
$\phi$:

$$
\overline{u}(\phi)
=
\phi\overline{u}.
$$

La funzione parametrica assume quindi la forma

$$
(\alpha,\beta,\phi)
\longmapsto
\texttt{solve\_alm}(\alpha,\beta,\phi).
$$

La struttura del solver rimane quella gi\`a utilizzata nel Caso Aula:
budget iniziale e bilanci temporali vengono rappresentati mediante
uguaglianze, i vincoli prudenziali mediante disuguaglianze e i limiti sul
funding mediante bound.

Il vincolo legacy richiede tuttavia particolare attenzione nella traduzione
computazionale. La condizione economica

$$
x_3+x_4\geq 35
$$

\noindent deve essere riscritta nella forma compatibile con

$$
A_{\mathrm{ub}}v\leq b_{\mathrm{ub}},
$$

\noindent ossia

$$
-x_3-x_4\leq -35.
$$

La trasformazione non modifica il vincolo economico, ma modifica il segno del
termine noto utilizzato dal solver.

Questa distinzione diventa rilevante quando si interpretano i marginali. Il
marginale restituito dal solver si riferisce al lato destro della
disuguaglianza codificata, cio\`e a $-35$. Se si vuole invece interpretare
l'effetto di una variazione del limite economico naturale

$$
L=35,
$$

\noindent nel vincolo

$$
x_3+x_4\geq L,
$$

\noindent occorre tener conto della relazione

$$
b_{\mathrm{ub}}=-L.
$$

Una variazione positiva di $L$ corrisponde quindi a una variazione negativa
del termine noto utilizzato dal solver. Il segno del marginale economico
rispetto a $L$ deve essere conseguentemente ricostruito a partire dalla
formulazione effettivamente codificata.

\subsection{Trasferire il metodo, non la soluzione}

Il passaggio dal Caso Aula al Caso TakeHome conserva la stessa sequenza
operativa:
\vspace{-0.3cm}
\begin{gather*}
\text{modello ALM}\\
\downarrow\\
\text{forma del solver}\\
\downarrow\\
\text{benchmark}\\
\downarrow\\
\text{funzione parametrica}\\
\downarrow\\
\text{sensitivit\`a e scenari}.
\end{gather*}


Non rimangono invece invariati i risultati, i vincoli attivi o il ruolo dei
singoli parametri.

Nel Caso TakeHome, il benchmark \`e determinato congiuntamente dal vincolo
prudenziale e dal vincolo legacy. Il funding viene utilizzato in una diversa
fase dell'orizzonte e la risposta ai due parametri di fabbisogno dipende dalla
loro collocazione temporale.

In particolare, uno shock terminale pu\`o essere assorbito inizialmente
utilizzando una giacenza disponibile alla data finale, senza richiedere una
modifica immediata dell'allocazione iniziale. Questo risultato non deve essere
interpretato come una ricomposizione successiva del portafoglio: le variabili
$x$ vengono scelte all'inizio dell'orizzonte e rimangono decisioni iniziali
del modello.

Il Caso TakeHome mostra quindi quali elementi possano essere trasferiti da un
problema all'altro e quali debbano essere ricostruiti. Rimangono invariati il
metodo di traduzione, la funzione del benchmark, la parametrizzazione e i
principi di verifica; cambiano la struttura economica del problema, i vincoli
rilevanti e il modo nel quale gli stress temporali si riflettono sulla
soluzione.


\section{Sintesi finale}

Questo capitolo ha reso operativo in Python il modello di programmazione
lineare e Asset--Liability Management sviluppato nei Capitoli 11 e 12.

Il passaggio centrale \`e consistito nella traduzione del problema economico
nella forma richiesta dal solver. Variabili decisionali, funzione obiettivo,
vincoli di uguaglianza, disuguaglianze e bound devono essere rappresentati in
modo coerente, mantenendo invariato il significato della formulazione
originaria.

La configurazione benchmark ha fornito il punto di riferimento per verificare
la correttezza della traduzione e interpretare la soluzione in termini di
allocazione dell'attivo, giacenze di liquidit\`a, ricorso al funding, vincoli
attivi e valori marginali.

L'introduzione di una funzione parametrica ha quindi permesso di passare dalla
soluzione di una singola configurazione allo studio di una famiglia di
problemi. Nel Caso Aula, i parametri $\theta$, $\phi$ e $\psi$ modificano
rispettivamente i fabbisogni di liquidit\`a, la capacit\`a di funding e il
costo del funding, intervenendo su componenti differenti della formulazione
del solver.

Su questa base, sono state distinte tre forme di analisi:

$$
\text{sensitivit\`a}
\qquad
\text{fattibilit\`a}
\qquad
\text{scenari}.
$$

L'analisi di sensitivit\`a mostra come cambiano il valore ottimo, il piano ALM
e la struttura dei vincoli attivi quando varia un parametro alla volta. La
frontiera di fattibilit\`a individua invece il massimo livello di stress
compatibile con l'esistenza di almeno una soluzione ammissibile. Gli scenari
consentono infine di studiare configurazioni nelle quali pi\`u driver vengono
modificati simultaneamente.

Il Caso TakeHome ha mostrato che la stessa architettura di analisi pu\`o
essere trasferita a un problema ALM differente, caratterizzato da una diversa
struttura dell'attivo, da fabbisogni collocati in date diverse e da ulteriori
vincoli strutturali.

Il percorso complessivo pu\`o essere sintetizzato come

$$
\begin{array}{c}
\text{modello ALM}
\\[3pt]
\downarrow
\\[3pt]
\text{forma del solver}
\\[3pt]
\downarrow
\\[3pt]
\text{benchmark verificato}
\\[3pt]
\downarrow
\\[3pt]
\text{funzione parametrica}
\\[3pt]
\downarrow
\\[3pt]
\text{sensitivit\`a, fattibilit\`a e scenari}.
\end{array}
$$

Il solver non sostituisce quindi il modello e non attribuisce autonomamente
significato economico ai risultati. La qualit\`a dell'analisi dipende dalla
corretta formulazione del problema, dalla coerenza della sua traduzione
computazionale e dalla verifica degli output prima della loro interpretazione.